% Auto-generado por scripts/extract-lessons-to-tex.ps1
% Fuente: HTML en carpeta L2-Teleologia-dosimetria-y-prescindencia-de-la-sancion-penal (sin modificar archivos originales).

\documentclass[11pt,a4paper]{article}
\usepackage[utf8]{inputenc}
\usepackage[T1]{fontenc}
\usepackage[spanish]{babel}
\usepackage{geometry}
\geometry{margin=2.5cm}
\usepackage{hyperref}
\hypersetup{colorlinks=true, linkcolor=blue, urlcolor=blue}

\title{\texttt{L2-Teleologia-dosimetria-y-prescindencia-de-la-sancion-penal}}
\author{Extracci\'on desde lecciones HTML}
\date{}

\begin{document}
\maketitle

\begin{small}
\noindent Este documento consolida el texto visible de los ocho componentes (01--08). Los formularios interactivos aparecen como texto; la l\'ogica JavaScript no se reproduce aqu\'i.
\end{small}
\bigskip
\section{Conceptos}
\noindent\textit{Archivo fuente:} \texttt{01-conceptos.html}

L2 — Conceptos\par





    L2 — Elemento 1 — Conceptos\par
    Teleología de la pena, dosimetría y prescindencia\par

    Introducción breve: cómo los fines del Código Penal y los principios de las sanciones condicionan la decisión sobre pena.\par


      Esta lección toma del Cap. 4 de Bernal el eje constitucional-criminal material que va del por qué se pune al\par
      cuánto y al si debe punirse. Los arts. 2 a 6 del C.P. expresan teleología legislativa; el juez opera\par
      dentro de mandatos de proporcionalidad, humanización y individualización sin sustituir la ley.\par



      1. Fines del código (art. 2 C.P.)\par


        El art. 2 orienta la interpretación sistemática del orden punitivo: protección de bienes jurídicos, respeto a derechos humanos y garantías,\par
        eficacia de la política criminal dentro del Estado social de derecho. No es un preámbulo ornamental: informa la ponderación\par
        en dosimetría y la revisión de penas desproporcionadas.\par


      Profundización: Bernal, Cap. 4, apartados sobre principios de las sanciones y función de la pena (aprox. pp. 285–320).\par




      2. Principios de las sanciones (art. 5 C.P.)\par



          Legalidad de la sanción\par

          Solo penas y medidas de seguridad previstas en ley; coherente con el bloque de constitucionalidad de L1.\par




          Proporcionalidad y resocialización\par

          La pena debe ser proporcional a la gravedad del hecho y la responsabilidad; el régimen debe favorecer la reintegración cuando corresponda.\par








      3. Funciones de la pena (art. 6 C.P.)\par



          FunciónNúcleoEfecto en litigio\par




          Prevención general positivaReforzar la confianza normativaArgumentos de política criminal en audiencias; no justifican pena ejemplar ilegítima\par

          Prevención general negativaDesincentivoDebe articularse con proporcionalidad estricta\par

          Prevención especialInhibición y, cuando aplica, tratamiento/resocializaciónIndividualización mediante circunstancias y medidas alternativas\par








      4. Dosimetría judicial\par


        La dosimetría es la operación de fijar la pena dentro del marco legal (mínimos, máximos, modalidades, concursos, circunstancias).\par
        Debe respetar: tipicidad ya declarada, proporcionalidad material, motivación suficiente y coherencia con el relato probatorio.\par



        Prescindencia / intervención mínima\par
        La idea de prescindencia o ultima ratio del derecho penal implica que solo debe acudirse a la pena cuando otras ramas\par
        del ordenamiento no bastan de modo razonable para la protección del bien jurídico. En defensa, articula pedidos de despenalización parcial,\par
        medidas sustitutivas o reconocimiento de proporcionalidad ex ante.\par





      5. Autocomprobación\par

      ¿Qué artículo del C.P. enumera explícitamente las funciones de la pena?\par

      Art. 5\par
      Art. 6\par





      Fuentes: Temario L2. Bernal, Cap. 4 (pp. 251–418; énfasis principios de sanciones, teleología y proporcionalidad, \textasciitilde{}285–341).\par
      C.P. arts. 2, 5 y 6.\par


    Siguiente →\par

\section{Explicaci\'on}
\noindent\textit{Archivo fuente:} \texttt{02-explicacion.html}

L2 — Explicación\par





    L2 — Elemento 2 — Explicación\par
    De la teleología a la sentencia: encadenamiento argumentativo\par

    Introducción breve: cómo se conectan fines constitucionales, arts. 2–6 C.P. y la motivación de la pena.\par


      El tribunal no “inventa” la teleología: aplica un sistema ya comprometido con dignidad humana y proporcionalidad (L1).\par
      La explicación dogmática ordena: (1) qué bien jurídico se protege; (2) qué gravedad fáctica y jurídica resulta probada; (3) qué margen legal\par
      de pena se abre; (4) cómo las circunstancias y la función preventiva especial ajustan el quantum dentro de ese margen.\par



      1. Proporcionalidad en dos niveles\par



          NivelPreguntaControl típico\par




          Proporcionalidad abstracta¿El tipo y el rango legal son compatibles con la gravedad media del delito?Control legislativo / inconstitucionalidad\par

          Proporcionalidad concreta¿La pena impuesta se ajusta a este hecho y este autor?Apelación, casación por error en el quantum\par








      2. Motivación de la pena\par


        La sentencia debe explicitar por qué se elige un punto dentro del marco legal: concurrencia de circunstancias, atenuantes, agravantes,\par
        delitos múltiples, reconocimiento de verdad, reparación. La ausencia de racionalidad dosimétrica expone el fallo a revocatoria\par
        por arbitrariedad o error judicial.\par





      3. Prescindencia operativa\par


        SubsidiariedadValorar si medidas no penales o penales mínimas satisfacen la finalidad constitucional.\par

        ProporcionalidadEvitar pena que no sea necesaria en sentido estricto para la prevención legítima.\par

        HumanizaciónCoherencia con art. 5 C.P. y estándar de trato digno (L1).\par






      4. Autocomprobación\par

      La proporcionalidad concreta se controla principalmente en:\par

      Primera instancia únicamente\par
      Segunda instancia y casación (error en el quantum)\par





      Fuentes: Bernal, Cap. 4 (principios de sanciones, función y límites de la pena). C.P. arts. 2, 5, 6.\par


    Siguiente →\par

\section{Exploraci\'on}
\noindent\textit{Archivo fuente:} \texttt{03-exploracion.html}

L2 — Exploración\par





    L2 — Elemento 3 — Exploración\par
    Dilema: “pena ejemplar” frente a proporcionalidad constitucional\par

    Escenario: delito mediático; presión social por penas máximas.\par


      El debate enfrenta la prevención general negativa entendida como intimidación colectiva frente al mandato de\par
      proporcionalidad y dignidad. La exploración entrena el argumento sin confundir política criminal con\par
      arbitrariedad judicial.\par



      Microcaso — Postura del sentenciador\par

       A) Imponer pena cercana al máximo para enviar “mensaje social”.\par
       B) Fijar pena según marco legal, gravedad concreta y motivación estricta de circunstancias.\par
       C) Minimizar siempre la pena sin análisis de gravedad.\par
      Ver retroalimentación\par





      Tabla de tensiones\par



          Argumento populistaContrapeso dogmático\par




          “La gente pide justicia” como único criterioLegalidad y motivación racional (arts. 5 y 6 C.P.; debido proceso)\par

          “Hay que castigar para educar a todos”Prevención general sin crueldad ni desproporción; función de la pena acotada\par








      Fuentes: Bernal, Cap. 4; C.P. arts. 2, 5, 6.\par


    Siguiente →\par

\section{Contexto}
\noindent\textit{Archivo fuente:} \texttt{04-contexto.html}

L2 — Contexto\par





    L2 — Elemento 4 — Contexto\par
    Política criminal, opinión pública y límites judiciales\par

    Contexto: cómo las reformas y el clima social presionan la dosimetría sin legitimar arbitrariedad.\par


      Las oleadas legislativas (mínimos obligatorios, figuras de especial gravedad, ampliación de marcos punitivos) modifican el\par
      “techo y piso” legal. El juez sigue obligado a motivar dentro del nuevo marco y a respetar principios constitucionales.\par
      La doctrina del Cap. 4 permite leer esas reformas en clave de compatibilidad con dignidad y proporcionalidad.\par



      Tipología de presiones\par



          FuenteEfecto en salaRespuesta técnica\par




          Medios y redesRiesgo de legitimar penas desproporcionadasMotivación escrita ajena a consignas; citar prueba y marco legal\par

          Discurso “cero tolerancia”Endurecimiento de práctica sentenciadorRevisión de constitucionalidad / casación si el marco o la motivación fallan\par

          Indicadores de gestión judicialIncentivos a sentencias rápidas o homogéneasIndividualización y derecho a la defensa no pueden sacrificarse\par








      Ejemplo sintético\par


        Una reforma sube el mínimo punitivo para un delito. El juez impone exactamente el mínimo sin analizar circunstancias de menor entidad.\par
        La defensa puede argumentar error en el quantum o déficit motivacional si existían bases para pena por debajo del tramo\par
        aplicado según el sistema legal vigente (concursos, atenuantes, etc.).\par





      Fuentes: Bernal, Cap. 4; temario L2.\par


    Siguiente →\par

\section{Aplicaci\'on}
\noindent\textit{Archivo fuente:} \texttt{05-aplicacion.html}

L2 — Aplicación\par





    L2 — Elemento 5 — Aplicación\par
    Clínica: alegatos de dosimetría en segunda instancia\par

    Plantilla para impugnar la pena por desproporción o déficit de motivación.\par


      Estructura su escrito en cuatro bloques: hechos probados relevantes para la gravedad; marco legal aplicable\par
      (mínimos, concursos, circunstancias); crítica a la lógica del sentenciador (saltos, omisiones, trato desigual de pruebas);\par
      petitorio concreto (subir, bajar o remitir para nueva dosimetría).\par



      Tesis modelo (adaptar)\par



-- La sentencia no explicitó la ponderación entre funciones de la pena del art. 6 C.P. y las circunstancias del arts. 59 y ss.\par


-- La pena impuesta se ubica en el tercio superior sin justificación suficiente frente a atenuantes acreditadas (enumerar).\par


-- Existe alternativa menos gravosa igualmente idónea para la prevención especial (fundamentar).\par





      Campos de trabajo\par



      Completitud mínima\par





      Fuentes: C.P. arts. 2, 5, 6; Bernal, Cap. 4.\par


    Siguiente →\par

\section{Prospectiva}
\noindent\textit{Archivo fuente:} \texttt{06-prospectiva.html}

L2 — Prospectiva\par





    L2 — Elemento 6 — Prospectiva\par
    Escenarios futuros de dosimetría y prescindencia\par


      Tendencias: mínimos punitivos sectoriales, uso de algoritmos de riesgo como insumo no vinculante (con debida\par
      prueba y contradicción), y litigio sobre trato digno en ejecución como continuación de la proporcionalidad de la pena impuesta.\par





          EscenarioImplicaciónEstrategia\par




          Más delitos con pena mínima altaMenos margen de individualizaciónLitigar circunstancias y concursos; control constitucional de tipos\par

          Expansión de medidas sustitutivasNuevo espacio de prescindenciaNegociación y verificación de requisitos legales\par

          Enfoque restaurativo reforzadoPosible convergencia con prevención especialArticular reparación con dosimetría\par








      Fuentes: Bernal, Cap. 4; temario L2.\par


    Siguiente →\par

\section{Ejercicios}
\noindent\textit{Archivo fuente:} \texttt{07-ejercicios.html}

L2 — Ejercicios\par





    L2 — Elemento 7 — Ejercicios\par
    Ejercicios con modelo\par

    Seis ítems sobre teleología, arts. 5–6 C.P. y dosimetría.\par


      1. Enumere tres funciones de la pena según art. 6 C.P.\par

      Ver modelo\par





      2. Diferencie proporcionalidad abstracta y concreta en una línea cada una.\par

      Ver modelo\par





      3. ¿Qué papel tiene el art. 2 C.P. en la dosimetría?\par

      Ver modelo\par





      4. Defina prescindencia penal en términos de política criminal.\par

      Ver modelo\par





      5. Indique dos vicios motivacionales que expongan la pena a revocatoria.\par

      Ver modelo\par





      6. Relacione art. 5 C.P. con estándar de dignidad humana (L1).\par

      Ver modelo\par




    Fuentes: C.P. arts. 2, 5, 6; Bernal, Cap. 4.\par

    Siguiente: Quiz →\par

\section{Quiz}
\noindent\textit{Archivo fuente:} \texttt{08-quiz.html}

L2 — Quiz\par





    L2 — Elemento 8 — Quiz\par
    Evaluación L2: teleología y dosimetría\par

    12 preguntas. Retroalimentación por ítem fallado.\par


      1. Los fines del código penal se encuentran principalmente en:\par

         A) Art. 2 C.P.\par
         B) Art. 5 C.P.\par
         C) Art. 6 C.P.\par


      2. Las funciones de la pena están desarrolladas en:\par

         A) Art. 5 C.P.\par
         B) Art. 6 C.P.\par
         C) Art. 29 C.P.\par


      3. Los principios de las sanciones (legalidad de la sanción, proporcionalidad, resocialización) están en:\par

         A) Art. 5 C.P.\par
         B) Art. 6 C.P.\par
         C) Art. 1 C.P.\par


      4. La proporcionalidad concreta se refiere a:\par

         A) Solo coherencia del tipo abstracto con el bien jurídico.\par
         B) Adecuación de la pena impuesta al hecho y autor probados.\par
         C) Solo decisiones del legislador.\par


      5. La prescindencia o lógica de última ratio implica:\par

         A) Usar siempre la pena más leve del código.\par
         B) Acudir al derecho penal solo cuando medios menos gravosos no basten razonablemente.\par
         C) Eliminar la prisión para todos los delitos.\par


      6. Una motivación de la pena vulnerable es:\par

         A) La que cita el marco legal y las circunstancias probadas.\par
         B) La que ignora atenuantes alegadas y acreditadas sin fundamento.\par
         C) La extensa en páginas aunque sea repetitiva.\par


      7. La prevención general negativa (art. 6) se orienta a:\par

         A) Solo tratamiento del condenado.\par
         B) Desincentivo general mediante la amenaza de la pena.\par
         C) Reparación integral de la víctima.\par


      8. Imponer pena cercana al máximo solo por presión mediática, sin fundamento en el caso:\par

         A) Es válido si el delito es notorio.\par
         B) Es crítico ante principios de proporcionalidad y debido proceso.\par
         C) Es obligatorio para restaurar la confianza ciudadana.\par


      9. La individualización judicial de la pena opera:\par

         A) Dentro del marco legal aplicable al caso.\par
         B) Sin límites si el juez considera justo.\par
         C) Solo en delitos políticos.\par


      10. La apelación respecto de la pena puede:\par

         A) Nunca modificar el quantum.\par
         B) Revisar el quantum y la motivación dosimétrica si hay error o arbitrariedad.\par
         C) Solo anular sin reenvío.\par


      11. El art. 5 C.P. dispone que el tratamiento del condenado debe favorecer su resocialización:\par

         A) Solo si el delito es leve.\par
         B) Cuando ello sea aplicable según la ley, en coherencia con la finalidad de la pena.\par
         C) Nunca, por ser contrario a la retribución.\par


      12. El art. 2 C.P. orienta la interpretación del código hacia:\par

         A) La máxima punición posible.\par
         B) Protección de bienes jurídicos, DDHH y eficacia racional en el Estado social de derecho.\par
         C) Excluir la participación de la víctima.\par


      Enviar\par


      Resultado\par









    Inicio L2\par

\end{document}

